\section{Variations with repetition}

In some ordered selections, the same object may be used more than once. Then the number of choices does not decrease from one position to the next.

\begin{definitionbox}
Let \(A\) be a finite set with \(n\) distinct elements. A variation with repetition of length \(k\) is an ordered \(k\)-tuple of elements of \(A\), where elements may repeat.
\end{definitionbox}

For example, if
\[
A=\{a,b\},
\]
then the variations with repetition of length \(3\) are
\[
aaa,\quad aab,\quad aba,\quad abb,\quad baa,\quad bab,\quad bba,\quad bbb.
\]
There are \(8\) of them.

\begin{theorembox}
The number of variations with repetition of length \(k\) chosen from \(n\) distinct elements is
\[
n^k.
\]
\end{theorembox}

There are \(n\) choices for the first position. Because repetition is allowed, there are again \(n\) choices for the second position, again \(n\) choices for the third position, and so on. Therefore the multiplication principle gives
\[
\underbrace{n\cdot n\cdots n}_{k\text{ factors}}=n^k.
\]

\begin{examplebox}
A three-symbol string is formed from the alphabet
\[
\{0,1,2,3,4,5,6,7,8,9\}.
\]
Symbols may repeat. There are
\[
10^3=1000
\]
such strings.
\end{examplebox}

\begin{remarkbox}
Use variations with repetition when order matters, there are \(k\) positions, and every position can be filled independently from the same set of \(n\) objects.
\end{remarkbox}
